%%%%%%%%%%%%%%%%%%%%%%%%%%%%%%%%%%%%%%%%%%%%%%%%%%%%%%%%%%%%%%%%%%%%%%%%
% Plantilla TFG/TFM
% Universidad de Murcia. Facultad de Informática
% Realizado por: José Manuel Requena Plens
% Modificado: Pablo José Rocamora Zamora
% Contacto: pablojoserocamora@gmail.com
%%%%%%%%%%%%%%%%%%%%%%%%%%%%%%%%%%%%%%%%%%%%%%%%%%%%%%%%%%%%%%%%%%%%%%%%

\chapter{Diseño y resolución del trabajo realizado}
\label{cap:diseno}

En este capítulo se expone la arquitectura técnica del sistema, el diseño del Grafo de Conocimiento y el modelado del comportamiento del agente inteligente responsable de traducir el lenguaje natural a consultas SPARQL.

\section{Arquitectura del Sistema}

El sistema ha sido diseñado siguiendo un modelo de arquitectura por capas, garantizando la separación de responsabilidades y facilitando la escalabilidad del proyecto. Los tres bloques principales son:

\begin{enumerate}
    \item \textbf{Capa de Presentación (Frontend):} Implementada con \textit{Streamlit}, actúa como el cliente de usuario, proveyendo la interfaz de chat interactiva.
    \item \textbf{Capa de Lógica de Negocio y Agente Inteligente:} Es el núcleo del sistema, orquestado mediante \textit{LangChain} en Python. Contiene la lógica de extracción de entidades, la integración con la API del LLM, el bucle de reflexión y la construcción del \textit{prompt} del sistema.
    \item \textbf{Capa de Datos y Grafo de Conocimiento:} Comprende el almacenamiento persistente RDF local (archivo \texttt{.ttl}) gestionado por la librería \textit{RDFLib}, así como las conexiones puntuales a la API de Wikidata para procesos de desambiguación.
\end{enumerate}

% ====================================================================
% NOTA PARA MARIO: 
% Aquí deberías insertar un Diagrama de Arquitectura de Software.
% ====================================================================

\begin{figure}[h]
    \centering
    % \includegraphics[width=0.8\textwidth]{ruta/a/tu/diagrama_arquitectura.png}
    \vspace{6cm} % Espacio temporal dejado en blanco
    \caption{[AÑADIR DIAGRAMA DE ARQUITECTURA AQUÍ]}
    \label{fig:arquitectura}
\end{figure}

\section{Diseño del Agente Autocorrector}

Uno de los principales desafíos técnicos del proyecto radica en la inestabilidad inherente de los LLMs a la hora de generar código con sintaxis estricta, como es el caso de SPARQL. Para solucionar esto, se ha diseñado un \textbf{flujo iterativo de reflexión y autocorrección}.

El flujo de control del agente sigue los siguientes pasos lógicos:
\begin{enumerate}
    \item Se recibe la pregunta en lenguaje natural.
    \item Se solicita al LLM que extraiga únicamente la entidad principal (ej. "Matrix").
    \item Se consulta Wikidata para obtener el ID de dicha entidad (ej. \texttt{wd:Q83495}).
    \item Se inyecta la pregunta, el ID de la entidad y el \textit{diccionario de propiedades} en un \textit{prompt} altamente restrictivo.
    \item El LLM genera una propuesta de consulta SPARQL.
    \item El sistema intenta ejecutar la consulta localmente en RDFLib.
    \item \textbf{Bucle de corrección:} Si la consulta falla (error de sintaxis o error de acceso a la ontología), el sistema no aborta la operación. Captura el \textit{traceback} del error, se lo envía de vuelta al LLM y le solicita que reescriba el SPARQL corrigiendo su propio fallo. Este proceso se repite un número máximo de iteraciones.
\end{enumerate}

% ====================================================================
% NOTA PARA MARIO: 
% Aquí deberías insertar un Diagrama de Actividad o de Secuencia
% mostrando el bucle del agente que acabo de describir.
% ====================================================================

\begin{figure}[h]
    \centering
    % \includegraphics[width=0.8\textwidth]{ruta/a/tu/diagrama_flujo.png}
    \vspace{6cm} % Espacio temporal dejado en blanco
    \caption{[AÑADIR DIAGRAMA DE FLUJO DEL AGENTE AQUÍ]}
    \label{fig:flujo_agente}
\end{figure}

\section{Diseño del Grafo de Conocimiento Local}

Para asegurar la privacidad, el rendimiento y la consistencia de las respuestas, el sistema no interroga directamente a Wikidata en su fase de resolución de consultas, sino que opera sobre un Grafo de Conocimiento local exportado previamente.

Este grafo local (\texttt{datos\_cine\_seguro.ttl}) actúa como una base de datos \textit{in-memory}. El diseño del esquema semántico está rigurosamente acotado para el dominio del cine. Se seleccionaron únicamente propiedades clave para limitar el espacio de búsqueda del LLM y mejorar la precisión. Entre estas propiedades mapeadas destacan:

\begin{itemize}
    \item \texttt{wdt:P57}: Director.
    \item \texttt{wdt:P577}: Fecha de publicación o estreno.
    \item \texttt{wdt:P161}: Miembro del reparto (actores).
    \item \texttt{wdt:P136}: Género cinematográfico.
    \item \texttt{wdt:P162}: Productor.
    \item \texttt{wdt:P345}: Identificador de IMDb.
\end{itemize}

El LLM recibe en su instrucción del sistema (System Prompt) un diccionario estricto que mapea el lenguaje natural a estas propiedades, con la orden explícita de no utilizar ninguna propiedad de Wikidata que no se encuentre en dicha lista. Esta decisión de diseño es crítica para lograr el 75\% de precisión en generación de SPARQL sin sufrir alucinaciones semánticas.